\section{背景とモデルの位置づけ}
\subsection{低圧プラズマの発光診断で求めるもの}
半導体のエッチング、成膜、チャンバー洗浄に用いる容量結合プラズマ（CCP）と誘導結合プラズマ（ICP）では、電子が電場からエネルギーを受け取り、原子・分子の励起、解離、電離を生じる。電子の特徴的なエネルギーはeV程度に達し、重粒子はそれより低い温度に保たれることが多い。この非平衡状態では、励起準位への衝突供給と放射・消光による損失を計算することが、発光を解釈する基本となる\cite{fujimoto,lieberman}。

OESは光を受動的に観測し、多数の原子線と分子帯を同時に取得する方法である。例えばF I線の強度には、F原子密度、高エネルギー電子による励起、親分子からの解離励起、衝突消光、視線内分布、波長応答が寄与する。分光診断で粒子密度や電子状態を求めるには、これらの寄与を反応と光学応答に分けて評価する\cite{coburn,fuller}。

OESCRはこの評価を、選択した発光種・準位と実験の観測条件に合わせて構成する。対象とする時間区間には、準位の定常近似を適用できる区間を選ぶ。幾何には、軸対称の殻モデルや外部で定めた投影行列を用いる。採用する励起・消光・捕獲の項は、放射寿命、衝突頻度、光学的厚さをもとに選定する。この手順によって、数Paの放電から、NF$_3$洗浄の例で参照する約133 Paの条件まで、それぞれの実験に即したモデルを組み立てられる\cite{an2023}。

\subsection{衝突放射モデルの考え方}
CRモデルは、励起準位への衝突供給、準位間移行、放射遷移、衝突失活を粒子収支で表す。放射損失が支配的な低密度極限では、基底状態からの直接励起と放射が釣り合うコロナモデルに近づく。密度が上がるにつれて、準位間衝突や準安定状態からの段階励起が重要になり、必要な状態と反応経路を収支へ加える\cite{fujimoto}。

OESCRでは、電子密度、EEDF、衝突相手の密度を与えたときの準位収支を線形方程式として解く。逆解析では、その外側で温度や密度を更新して速度係数と収支を再評価する。この構成により、一回の準位計算の見通しを保ちながら、非線形な発光応答に基づく状態推定を行える。

\begin{table}[htbp]
\centering\small
\begin{tabularx}{\linewidth}{p{26mm}YYY}
\toprule
モデル & 主に解く量 & 主な用途 & OESCRとの接点 \\
\midrule
LTE・Boltzmann分布 & 温度で定まる準位占有率 & 熱平衡に近い発光の解析 & 非平衡収支と比較する基準 \\
コロナモデル & 直接励起と放射の釣合い & 放射損失が支配する希薄系の線強度 & CRの採用項を限定した極限 \\
詳細CRモデル & 多数準位・多価数の占有率 & 機構を詳しく含めた分光診断 & 対象準位を選ぶ縮約の出発点 \\
0次元化学モデル & 化学種密度、粒子・電力収支 & 組成、解離率、運転条件の評価 & 供給粒子密度や電子状態の入力源 \\
流体・輸送モデル & 密度、流束、場の空間分布 & 放電構造と装置条件の評価 & 密度・温度分布を殻状態へ、観測配置を投影重みへ反映 \\
PIC/MCC & 荷電粒子の運動、衝突、電場 & 非局所加熱、シース、電子動力学 & 電子分布等を発光計算へ入力 \\
Boltzmann（BOLSIG+） & 二項近似の電子速度分布 & 指定した均一場・$E/N$下の輸送・反応係数 & 表形式EEDFと率積分の比較 \\
OESCR & 準位密度、測定スペクトル、指定した推定量 & 分光解釈、診断設計、機構・校正の感度比較 & 物理状態と観測値を結ぶ計算 \\
\bottomrule
\end{tabularx}
\caption{各モデルが担当する問題と、OESCRへの入出力の接点\cite{fujimoto,lieberman,hagelaar,arellano}。外部モデルの結果は、対応する密度・EEDF・幾何をcaseへ与えて利用する。}
\label{tab:models}
\end{table}

\subsection{利用目的に応じた計算の組み立て}
発光機構を調べる場合は、基底励起、準安定からの励起、消光等を選んで順解析し、準位収支と線強度の変化を比較する。測定から状態量を求める場合は、既知の供給密度と校正条件を固定し、観測が支える温度・密度を推定対象にする。診断設計では、異なる励起感度を持つ線、参照種、視線を加えたときのJacobianやprofileの変化を調べる。

複数ガスへの展開も同じ考え方に従う。各ガスの状態・反応データを用意し、共通の電子状態、混合相手による消光、観測応答を設定することで、同一の解析手順を使える。対象条件の選定とデータの根拠を明示することが、計算結果を実験の解釈へつなぐ要点となる。
